% GENERATED FILE. Edit canonical lessons, then run npm run generate:books.
% canonical-source-hash: fnv1a64:310bd3fb5247de3b
% canonical-lessons: TE-C142-kastapade, TE-C142-telivaina, TE-C142-bhayankaramaina, TE-C142-nallani, TE-C142-tellani

\chapter{Hard-working, Clever, Terrible, Black, White}
\label{ch:te-hard-working-clever-terrible-black-white}
\hlchaptermodality{142}

\begin{chapteropening}
\textbf{By the end of this chapter:} \emph{I can say hard-working, clever, terrible, black and white.}

Complete the last lesson of chapter 142: I can say hard-working, clever, terrible, black and white.
\end{chapteropening}

\section[\texorpdfstring{\te{కష్టపడే} (kaṣṭapaḍē)}{కష్టపడే (kaṣṭapaḍē)}]{\te{కష్టపడే} (kaṣṭapaḍē) — hard-working}

\label{lesson:TE-C142-kastapade}



\begin{quote}
Before the new one: say the Telugu for courage; brave, then the Telugu for lazy.
\end{quote}

\subsection*{You'll want to know: \te{కష్టపడే}}
\textbf{\te{కష్టపడే}} — \emph{kaṣṭapaḍē} — \textquotedblleft{}hard-working\textquotedblright{}.

\begin{cousinweb}
\textbf{\te{కష్టపడే}} (\emph{kaṣṭapaḍē}) is Sanskrit \emph{kaṣṭa}, \textquotedblleft{}hardship\textquotedblright{} (CDIAL 2978), with native \emph{paḍu}, \textquotedblleft{}to suffer, undergo\textquotedblright{} (DEDR 3853): \textquotedblleft{}one who undergoes hardship\textquotedblright{}.
\end{cousinweb}

\begin{grammarlens}[title={What you've built}]
One more word for this chapter.
\end{grammarlens}

\subsection*{Your turn}
\begin{itemize}
\raggedright
  \item \emph{Say it:} \emph{kaṣṭapaḍē}
  \item \emph{Say it:} \emph{kaṣṭapaḍē}, once more
  \item \emph{Say it:} say \emph{kaṣṭapaḍē}
  \item \emph{Recall:} say the Telugu for courage; brave, then the Telugu for lazy, then say \emph{kaṣṭapaḍē} again
\end{itemize}

\begin{tcolorbox}[breakable,colback=teal!4,colframe=teal!35!black,title={Before you move on}]
What does \te{కష్టపడే} mean? (\textquotedblleft{}Hard-working\textquotedblright{}.) Say it once more.
\end{tcolorbox}
\section[\texorpdfstring{\te{తెలివైన} (telivaina)}{తెలివైన (telivaina)}]{\te{తెలివైన} (telivaina) — clever}

\label{lesson:TE-C142-telivaina}



\begin{quote}
Before the new one: say the Telugu for lazy, then the Telugu for hard-working.

\emph{Read it:} \textbf{\te{ఇది} \te{తలుపు}. \te{దీపం} \te{ఇక్కడ} \te{ఉంది}.} aloud as one run
\end{quote}

\subsection*{You'll want to know: \te{తెలివైన}}
\textbf{\te{తెలివైన}} — \emph{telivaina} — \textquotedblleft{}clever\textquotedblright{}.

\begin{cousinweb}
\textbf{\te{తెలివైన}} (\emph{telivaina}) is native Dravidian: \emph{telivi}, \textquotedblleft{}understanding, intelligence\textquotedblright{}, which DEDR 3433 lists under \emph{teliyu}, \textquotedblleft{}to be known\textquotedblright{}, and \emph{ayina}, \textquotedblleft{}which is\textquotedblright{} (DEDR 333).
\end{cousinweb}

\begin{grammarlens}[title={What you've built}]
One more word for this chapter.
\end{grammarlens}

\subsection*{Your turn}
\begin{itemize}
\raggedright
  \item \emph{Say it:} \emph{telivaina}
  \item \emph{Say it:} \emph{telivaina}, once more
  \item \emph{Say it:} say \emph{telivaina}
  \item \emph{Recall:} say the Telugu for lazy, then the Telugu for hard-working, then say \emph{telivaina} again
\end{itemize}

\begin{tcolorbox}[breakable,colback=teal!4,colframe=teal!35!black,title={Before you move on}]
What does \te{తెలివైన} mean? (\textquotedblleft{}Clever\textquotedblright{}.) Say it once more.
\end{tcolorbox}
\section[\texorpdfstring{\te{భయంకరమైన} (bhayaṅkaramaina)}{భయంకరమైన (bhayaṅkaramaina)}]{\te{భయంకరమైన} (bhayaṅkaramaina) — terrible, dreadful}

\label{lesson:TE-C142-bhayankaramaina}



\begin{quote}
Before the new one: say the Telugu for hard-working, then the Telugu for clever.
\end{quote}

\subsection*{You'll want to know: \te{భయంకరమైన}}
\textbf{\te{భయంకరమైన}} — \emph{bhayaṅkaramaina} — \textquotedblleft{}terrible, dreadful\textquotedblright{}.

\begin{cousinweb}
\textbf{\te{భయంకరమైన}} (\emph{bhayaṅkaramaina}) is \emph{bhayaṅkara}, Sanskrit for \textquotedblleft{}causing fear\textquotedblright{}, built on Sanskrit \emph{bhaya}, \textquotedblleft{}fear\textquotedblright{} (CDIAL 9390), with native \emph{ayina}, \textquotedblleft{}which is\textquotedblright{} (DEDR 333).
\end{cousinweb}

\begin{grammarlens}[title={What you've built}]
One more word for this chapter.
\end{grammarlens}

\subsection*{Your turn}
\begin{itemize}
\raggedright
  \item \emph{Say it:} \emph{bhayaṅkaramaina}
  \item \emph{Say it:} \emph{bhayaṅkaramaina}, once more
  \item \emph{Say it:} say \emph{bhayaṅkaramaina}
  \item \emph{Recall:} say the Telugu for hard-working, then the Telugu for clever, then say \emph{bhayaṅkaramaina} again
\end{itemize}

\begin{tcolorbox}[breakable,colback=teal!4,colframe=teal!35!black,title={Before you move on}]
What does \te{భయంకరమైన} mean? (\textquotedblleft{}Terrible, dreadful\textquotedblright{}.) Say it once more.
\end{tcolorbox}
\section[\texorpdfstring{\te{నల్లని} (nallani)}{నల్లని (nallani)}]{\te{నల్లని} (nallani) — black, dark}

\label{lesson:TE-C142-nallani}



\begin{quote}
Before the new one: say the Telugu for clever, then the Telugu for terrible.
\end{quote}

\subsection*{You'll want to know: \te{నల్లని}}
\textbf{\te{నల్లని}} — \emph{nallani} — \textquotedblleft{}black, dark\textquotedblright{}.

\begin{cousinweb}
\textbf{\te{నల్లని}} (\emph{nallani}) is native Dravidian, one word with Tamil \emph{nallam} (DEDR 3613). DEDR lists Telugu \emph{nalla} and \emph{nallani}, \textquotedblleft{}black\textquotedblright{}.
\end{cousinweb}

\begin{grammarlens}[title={What you've built}]
One more word for this chapter.
\end{grammarlens}

\subsection*{Your turn}
\begin{itemize}
\raggedright
  \item \emph{Say it:} \emph{nallani}
  \item \emph{Say it:} \emph{nallani}, once more
  \item \emph{Say it:} say \emph{nallani}
  \item \emph{Recall:} say the Telugu for clever, then the Telugu for terrible, then say \emph{nallani} again
\end{itemize}

\begin{tcolorbox}[breakable,colback=teal!4,colframe=teal!35!black,title={Before you move on}]
What does \te{నల్లని} mean? (\textquotedblleft{}Black, dark\textquotedblright{}.) Say it once more.
\end{tcolorbox}
\section[\texorpdfstring{\te{తెల్లని} (tellani)}{తెల్లని (tellani)}]{\te{తెల్లని} (tellani) — white, bright}

\label{lesson:TE-C142-tellani}



\begin{quote}
Before the new one: say the Telugu for terrible, then the Telugu for black.
\end{quote}

\subsection*{You'll want to know: \te{తెల్లని}}
\textbf{\te{తెల్లని}} — \emph{tellani} — \textquotedblleft{}white, bright\textquotedblright{}.

\begin{cousinweb}
\textbf{\te{తెల్లని}} (\emph{tellani}) is native Dravidian. DEDR 3433 lists \emph{tella} and \emph{tellani}, \textquotedblleft{}white\textquotedblright{}, under \emph{teli}, \textquotedblleft{}white, pure\textquotedblright{}, with Tamil \emph{teḷ}, \textquotedblleft{}to become clear\textquotedblright{}.
\end{cousinweb}

\begin{grammarlens}[title={What you've built}]
One more word for this chapter.
\end{grammarlens}

\subsection*{Your turn}
\begin{itemize}
\raggedright
  \item \emph{Say it:} \emph{tellani}
  \item \emph{Say it:} \emph{tellani}, once more
  \item \emph{Say it:} say \emph{tellani}
  \item \emph{Recall:} say the Telugu for terrible, then the Telugu for black, then say \emph{tellani} again
\end{itemize}

\begin{tcolorbox}[breakable,colback=teal!4,colframe=teal!35!black,title={Before you move on}]
What does \te{తెల్లని} mean? (\textquotedblleft{}White, bright\textquotedblright{}.) Say it once more.
\end{tcolorbox}
